\chapter{Machine Learning para avaliação de jogadores e alinhamentos}
\label{chap:desempenho}

O módulo de avaliação de jogadores e alinhamentos de AAUTIA implementa duas tarefas distintas, documentadas em
\path{docs/PERFORMANCE_EVALUATION.md}: previsão da produção individual
e avaliação da associação dos cinco iniciais com a margem final do jogo.
Não são reconstruídos todos os alinhamentos utilizados durante a partida.

\section{Previsão da produção individual}

Cada observação corresponde à participação de um jogador num jogo e
exige pelo menos cinco participações anteriores na mesma temporada.
Os alvos são pontos, ressaltos totais e assistências, tratados separadamente
para evitar impor pesos arbitrários numa única classificação de jogador.
A previsão é condicional à participação: não estima convocatória,
disponibilidade ou lesões.

As características incluem jogar em casa, dias de descanso, número de
participações históricas limitado a dez e médias das últimas cinco e dez
participações para pontos, ressaltos, assistências, roubos de bola,
desarmes de lançamento, perdas de bola e tentativas de lançamento de campo.
O histórico acompanha o jogador nas mudanças de equipa dentro da mesma
temporada e reinicia na época seguinte. Os minutos não são usados, pois
a codificação desse campo na fonte necessita de validação adicional.
Resumos de carreira e estatísticas da partida prevista também são excluídos.

São comparadas regressão Ridge, Random Forest e média das últimas dez
participações. A seleção minimiza a média dos erros absolutos médios
(MAE) dos três alvos na validação. Esta média usa as unidades originais,
pelo que os pontos podem ter maior influência no critério. A referência
histórica pode ser selecionada caso supere os modelos aprendidos.
A imputação e a normalização aplicável são ajustadas dentro do processo
de treino, seguindo o protocolo temporal comum.

O resumo por jogador apresenta produção observada e resíduos, incluindo
pontos observados menos previstos. Um resíduo positivo indica produção
superior à expectativa do modelo, sem demonstrar qualidade intrínseca ou
contributo causal para a vitória. A comparação agregada dispõe de uma
marca de elegibilidade para jogadores com pelo menos 20 participações
avaliadas. A exclusão do início de época e de jogadores com pouco histórico
limita a representatividade.

\section{Avaliação dos cinco iniciais}

Uma partida só é aceite quando existem exatamente cinco titulares distintos
por equipa, sem jogadores partilhados entre os dois lados. Os jogos com
registos inválidos ou incompletos são excluídos e identificados num ficheiro
próprio. Existe uma lacuna importante nos titulares de 2021--22, comum às
fontes individuais básica e estendida.

O alvo é a margem final de pontos da equipa da casa. A representação
aditiva atribui \(+1\) a cada titular da casa e \(-1\) a cada titular
visitante, originando o modelo
\[
\widehat{d}_g=b+\sum_{i\in H_g}\beta_i-\sum_{i\in A_g}\beta_i,
\]
em que \(H_g\) e \(A_g\) são os cinco iniciais e \(b\) representa a vantagem
de casa ajustada. Uma segunda representação acrescenta indicadores das
combinações completas de cinco jogadores. A regressão Ridge regulariza
os coeficientes de jogadores e combinações pouco observados.

A representação e a regularização são selecionadas pelo menor RMSE na
validação, incluindo a margem média em casa como referência elegível.
Jogadores ou combinações desconhecidos recebem contribuição aprendida
nula, sendo contabilizados nos resultados. O resumo de cada quinteto soma
os coeficientes aplicáveis sem o intercepto e apresenta também a margem
final observada nos jogos em que esse quinteto iniciou a partida.
A marca de elegibilidade exige pelo menos dez jogos de teste como titular
para a combinação.

Estes coeficientes são associações com o resultado do jogo completo,
não avaliações por posse ou estimativas de desempenho conjunto durante
o tempo efetivamente passado em campo. Banco, adversários e decisões do
treinador influenciam a margem. Os titulares são conhecidos retrospetivamente,
sem verificação da sua disponibilidade antes do início do jogo.

\section{Interfaces para o planeamento}

O procedimento \path{scripts/train_performance.py} guarda modelos,
metadados, previsões, coeficientes, cobertura e exclusões. O comando
\path{scripts/evaluate_lineup.py} aceita cinco identificadores por lado
e devolve a margem esperada do jogo e os jogadores desconhecidos; não
valida a pertença ao plantel nem fornece uma probabilidade de vitória.

O módulo de planeamento de rotações reutiliza as previsões individuais e os coeficientes aditivos
dos titulares. A implementação rejeita modelos com interações de quintetos
em vez de omitir esses termos. A passagem das previsões para uma utilidade
de minutos é uma hipótese do planeamento, descrita no capítulo seguinte,
e não uma propriedade aprendida pelo modelo de produção.
